\input{preamble}

\begin{document}

\title{Commutative algebra exam}
\maketitle

\section{Irreducible polynomials and Galois
theory}

\begin{exercise}
\label{exercise-rational-functions-irreducible}
Let $k=\mathbb{C}(t)$ be the field of
rational functions on $\mathbb{C}$.
Prove that for any $n$ there exists
an irreducible polynomial
$P(t) \in k[t]$ of degree $n$.
\end{exercise}

\begin{proof}
Let $P(x) \in k[x]=\mathbb{C}(t)[x]$
be $P(x) = x^n-t$.
This polynomial has no root in
$k=\mathbb{C}(t)$,
for if $(f/g)^n=t$ then 
$n(\operatorname{deg}f-
\operatorname{deg}g)=1$,
which is impossible for $n>1$.
But unfortunately this doesn't prove
that it's irreducible: in $\mathbb{R}[x]$,
the polynomial
$(x^2 + 1)(x^2+2)=x^4+3x^2+2$
is reducible and has no roots.

To prove $x^n + t$ is irreducible we shall
use the Gauss lemma. Forget the above
notation and suppose that
$fg = x^n + t$, where
$f,g \in \mathbb{C}(t)[x]$.
We want to clear the denominators in
the coefficients of $f$ and $g$,
so we multiply by the product $D$ of all
such denominators and obtain
$$
f_1g_1 = (x^n + t)D,
$$
where $D \in \mathbb{C}[t]$ and
$f_1,g_1 \in \mathbb{C}[t][x]$.
Now we apply Gauss lemma, that is, the
content of the product is the product
of the contents. Since $x^n +t$ is primitive
(i.e. has content $1$), then the
content of $f_1g_1$ equals the content
of $D$, but $D$ is a constant
(constants here are polynomials in
$\mathbb{C}[t]$).
We obtain
$$
FG = x^n + t
$$
where $F$ and $G$ are the primitive parts
of $f_1$ and $g_1$. This holds up to
multiplication by a complex number.

Now we reduce modulo $t$ and obtain
$$
[F][G] = [x^n],
$$
and conclude that, as classes,
$$
[F]=[cx^m],\qquad 
[G]=[c^{-1}x^{n-m}].
$$
Thus we see that each of $F,G$
lost their constant term when we reduced,
i.e. that their constant terms are divisible
by $t$. This means that $t^2$ divides
the constant term of the product $FG$.
But this is impossible, since the constant
term of $FG$ is $t$. 

The only way
this makes sense is if $m=0$ or $m=n$.
To conclude we note that reduction modulo
$t$ preserves the degree of both $F$, and 
$G$ because their leading coefficients
multiply to $1$. Thus $m=0$ or $m=n$
imply that $F$ or $G$ is a unit, which
in turn implies so for $f$ or $g$.
\end{proof}

\begin{exercise}
\label{exercise-smooth-irreducible-plane}
Prove that for any $d \in  \mathbb{Z}^{>0}$
there exists an irreducible polynomial
$P \in \mathbb{C}[x,y]$ such that
$P(x,y)=0$ defines a smooth submanifold
in $\mathbb{C}^2$.
\end{exercise}

\begin{proof}
$x^n+y$ since $\partial_y =1$.
This is an irreducible polynomial:
if $x^n+y = fg$, then either of $f,g$
must have no $y$ variable, say $f \in
\mathbb{C}[x]$,
and the other one must be degree 1 in $y$.
Suppose $g = y f' + f''$ for $f',f'' \in
\mathbb{C}[x]$. But then
$x^n+y = fg = yff' + f f''$,
which means that $f$ cannot have positive
degree in $x$, but $f \in \mathbb{C}[x]$,
so it is a unit.
\end{proof}

\begin{exercise}
\label{exercise-smooth-irreducible-plane}
Prove that for any $d \in  \mathbb{Z}^{>0}$
there exists an irreducible polynomial
$P \in \mathbb{C}[x,y]$ such that
$P(x,y)=0$ defines a smooth submanifold
in $\mathbb{C}^2$.
\end{exercise}

\begin{proof}
$x^n+y$ since $\partial_y =1$.
This is an irreducible polynomial
because $y$ is irreducible
More precisely,
if $x^n+y = fg$, then either of $f,g$
must have no $y$ variable, say $f \in
\mathbb{C}[x]$,
and the other one must be degree 1 in $y$.
Suppose $g = y g' + g''$ for $g',g'' \in
\mathbb{C}[x]$. But then
$x^n+y = fg = yfg' + f g''$,
which means that $f$ cannot have positive
degree in $x$, but $f \in \mathbb{C}[x]$,
so it is a unit.
\end{proof}


\end{document}
